\title{Controlling Text-to-Image Diffusion by Orthogonal Finetuning}

\begin{abstract}
Large-scale text-to-image diffusion models have demonstrated remarkable proficiency in synthesizing photorealistic visuals from textual prompts. However, efficiently directing or adapting these advanced models to fulfill diverse downstream objectives remains a significant unresolved question. In response, we propose a systematic finetuning strategy termed Orthogonal Finetuning~(OFT) for the adaptation of text-to-image diffusion models in downstream scenarios. Unlike prevailing approaches, OFT can theoretically maintain hyperspherical energy, reflecting the relationships between neurons on the unit hypersphere. Our analysis suggests that preserving this attribute is vital for retaining the semantic generative capabilities of text-to-image diffusion frameworks. To further enhance the stability of the finetuning process, we introduce Constrained Orthogonal Finetuning (COFT), which adds an explicit constraint on the radius of the hypersphere. We focus on two key finetuning applications: subject-driven generation, which involves producing images of a particular subject in varied contexts given a limited number of reference images and a text prompt, and controllable generation, where the model integrates supplementary control signals. Experimental results indicate that our OFT approach achieves superior generative fidelity and faster convergence relative to competing methods.
\end{abstract}

\section{Introduction}

State-of-the-art text-to-image diffusion models~\cite{saharia2022photorealistic,ramesh2022hierarchical,rombach2022high} have achieved outstanding results in leveraging textual instructions for generating high-quality, detailed images. Nonetheless, textual prompts alone often lack the specificity required to afford detailed and precise manipulations over generated outputs. To mitigate this issue, our work addresses two main categories of text-to-image generation tasks:

\begin{itemize}[leftmargin=*,nosep]

    \item \textbf{Subject-driven generation}~\cite{ruiz2023dreambooth}: Given a small set of images depicting a subject, the objective is to synthesize new images representing the same subject situated in alternative contexts, utilizing an accompanying textual description as guidance.
    \item \textbf{Controllable generation}~\cite{zhang2023adding,mou2023t2i}: When provided with an extra control input (such as canny edge maps or segmentation masks), the aim is to direct the model to produce images conforming to this control information in addition to the text prompt.
\end{itemize}

At their core, both tasks focus on the challenge of finetuning text-to-image diffusion models in a way that maintains the generative strengths achieved during pretraining. The primary requirements for an effective finetuning approach can be outlined as follows: (1) \emph{training efficiency}, meaning a reduction in the number of trainable parameters and training epochs needed, and (2) \emph{generalizability preservation}, which refers to retaining the model’s ability to produce high-quality and varied outputs. Typical strategies for finetuning involve either incrementally updating neuron weights with a low learning rate (\emph{e.g.}, \cite{ruiz2023dreambooth}) or introducing a lightweight module that utilizes re-parameterized weights (\emph{e.g.}, \cite{hulora2022,zhang2023adding}). Despite being straightforward, these methods fall short in reliably safeguarding the generative capabilities acquired during pretraining. Moreover, there is a notable absence of a principled framework for devising robust finetuning techniques or selecting appropriate hyperparameters, such as optimal epoch counts. A central obstacle lies in the absence of quantitative metrics that assess how well finetuning conserves the pretrained generative properties. Most prevailing methods implicitly posit that a smaller Euclidean distance between finetuned and pretrained models correlates with superior retention of prior abilities; as a consequence, finetuning is generally conducted with extremely low learning rates. However, this assumption, while sometimes valid, does not comprehensively reflect the extent of semantic preservation. We contend that Euclidean proximity to the original model does not suffice as an indicator of semantic retention. Employing a more structurally informed metric to distinguish between finetuned and pretrained models could significantly enhance both the preservation of pretrained performance and the stability of the finetuning process.

Drawing on empirical findings that hyperspherical similarity effectively captures semantic relationships~\cite{liu2017deep,liu2018decoupled,chen2020angular}, we utilize hyperspherical energy~\cite{liu2018learning} to analyze the pairwise associations between neurons. Specifically, hyperspherical energy aggregates the hyperspherical similarity (such as cosine similarity) among all neuron pairs within a layer, thereby quantifying the degree of neuron uniformity on the unit hypersphere~\cite{liu2021learning}. We posit that optimal finetuning should result in minimal change to the hyperspherical energy relative to the pretrained model. One straightforward strategy would be to introduce a regularization term enforcing constant hyperspherical energy during finetuning; however, such an approach does not necessarily ensure the energy difference is adequately controlled. To address this, we leverage a key invariance property: applying a uniform orthogonal transformation to all neurons preserves the pairwise hyperspherical similarity. Based on this invariance, we introduce Orthogonal Finetuning~(OFT), a method for adapting large text-to-image diffusion models to new tasks without modifying their hyperspherical energy. The core principle is to learn an orthogonal transformation, shared across a layer, that maintains the angular relationships among neurons. OFT may also be interpreted as realigning the canonical coordinate system of neurons within a layer. Furthermore, recognizing that maintaining a small Euclidean distance between finetuned and pretrained models tends to safeguard pretraining advantages, we propose Constrained Orthogonal Finetuning~(COFT): a variant of OFT that restricts the finetuned model within a hyperspherical region of fixed radius centered at the pretrained neurons.

The rationale underlying the use of orthogonal transformations for finetuning neurons draws, in part, from insights offered by the 2D Fourier transform. Through this decomposition, an image can be represented by both its magnitude and phase spectra. Crucially, the phase spectrum, which encodes angular relationships between inputs and the basis, retains the essential semantic information. As evidence, reconstructing an image using its phase spectrum paired with an arbitrary magnitude spectrum still preserves its semantics. This observation implies that the orientation of neurons is pivotal for modulating the semantics in generated images—a key objective in subject-driven and controllable image generation tasks. Nonetheless, arbitrary alteration of neuron orientations can undermine the model’s generative performance that was established during pretraining. To avoid excessive changes, we advocate for preserving the pairwise angles between neurons. This principle is motivated by the hypothesis that such angular relationships are vital for encoding the knowledge within neural networks. Building on this foundation, we propose to apply a uniform orthogonal transformation to neurons within each layer, thereby ensuring the conservation of hyperspherical energy.

Further motivation is drawn from orthogonal over-parameterized training~\cite{liu2021orthogonal}, a technique in which classification networks are trained from scratch by orthogonally transforming an initially random neural network. Random initialization is theoretically associated with low expected hyperspherical energy, and \cite{liu2021orthogonal} aims to maintain this low energy state throughout training, a property correlated with improved generalization in classification tasks~\cite{liu2018learning,lin2020regularizing}. Their findings demonstrate that orthogonal transformations possess sufficient expressive power to optimize neural networks for classification. In contrast, our work is oriented towards finetuning text-to-image diffusion models to enable enhanced control and superior generative capability for downstream tasks. We would like to highlight two fundamental distinctions between OFT and the approach in \cite{liu2021orthogonal}. First, while \cite{liu2021orthogonal} is dedicated to reducing hyperspherical energy, OFT instead strives to maintain the hyperspherical energy established by the pretrained model to avoid disrupting its learned structures—since reducing this energy during diffusion model finetuning may compromise their intrinsic semantic organization. Second, OFT incorporates a constraint to limit the divergence from the original pretrained model, resulting in a constrained formulation, in contrast to the unconstrained nature of \cite{liu2021orthogonal}. Finetuning success thus hinges on appropriately balancing flexibility with stability; we contend that our OFT methodology accomplishes this balance effectively. The specific contributions of our work are summarized below:

\begin{itemize}[leftmargin=*,nosep]

    \item We introduce Orthogonal Finetuning, a new approach for enhancing controllability in text-to-image diffusion models. Additionally, we present a constrained version that enforces a bound on the angular shift relative to the original pretrained model to boost finetuning stability.
    \item In contrast to conventional finetuning strategies, OFT fine-tunes models while rigorously maintaining the hyperspherical energy—a property we empirically demonstrate to be crucial for retaining the semantic generative capacity of the base model.
    \item Our experiments deploy OFT across two domains: subject-driven generation and controllable generation. Through extensive evaluation, we establish that OFT leads to marked gains in image quality, training speed, and the robustness of finetuning compared to existing techniques. Notably, OFT also exhibits greater sample efficiency, requiring fewer training images and epochs to achieve strong convergence.
    \item For the task of controllable generation, we devise a control consistency metric, which quantitatively assesses controllability by extracting the control signal from synthesized images and juxtaposing it against the original input. Our metric substantiates the superior controllability afforded by OFT.
\end{itemize}

\section{Related Work}

\textbf{Text-to-image diffusion models}. Substantial advances~\cite{saharia2022photorealistic,ramesh2022hierarchical,rombach2022high,gu2022vector,nichol2022glide} have been achieved in the area of text-to-image synthesis, largely driven by the evolution of diffusion-based generative techniques~\cite{ho2020denoising,song2021score,song2021denoising,dhariwal2021diffusion} in conjunction with progress in vision-language pretraining~\cite{su2020vlbert,lu2019vilbert,radford2021learning,wang2021simvlm,li2022blip,li2023blip,alayrac2022flamingo,schuhmann2022laion}. Models such as GLIDE~\cite{nichol2022glide} and Imagen~\cite{saharia2022photorealistic} implement diffusion mechanisms directly in pixel space. Specifically, GLIDE employs joint training of the text encoder and a diffusion prior using text-image pairs, whereas Imagen leverages a fixed pretrained text encoder. Meanwhile, architectures like Stable Diffusion~\cite{rombach2022high} and DALL-E2~\cite{ramesh2022hierarchical} focus on diffusion processes within a latent space. Stable Diffusion utilizes VQ-GAN~\cite{esser2021taming} to construct its latent representation, and DALL-E2 relies on CLIP~\cite{radford2021learning} to generate a multimodal embedding that encodes both images and text. Beyond diffusion, generative adversarial networks~\cite{reed2016generative,zhang2017stackgan,xu2018attngan,li2019controllable} and autoregressive approaches~\cite{ramesh2021zero,ding2021cogview,wu2022nuwa,yuscaling} have been explored for text-conditioned image generation. OFT operates as a model-agnostic finetuning scheme, suitable for application to any text-to-image diffusion model.

\textbf{Subject-driven generation}. To avoid alterations to the primary subject, some works~\cite{avrahami2022blended,nichol2022glide} introduce user-provided masks as auxiliary conditioning. Inversion techniques~\cite{choi2021ilvr,dhariwal2021diffusion,ramesh2022hierarchical,gal2022image} are employed to adjust context while keeping the subject largely intact, and \cite{hertz2022prompt} achieve both local and global editing without requiring explicit masks. However, these methods generally fail to preserve subject identity details faithfully. Pivotal Tuning~\cite{roich2022pivotal} adapts the generator within the neighborhood of an inverted latent code, incorporating extra regularization to maintain identity information. Likewise, \cite{nitzan2022mystyle} crafts personalized generative priors by learning from a set of an individual's facial images. Instance-level diversity is explored in \cite{casanova2021instance}, although crucial subject attributes might be compromised. DreamBooth~\cite{ruiz2023dreambooth} advances personalization by leveraging a custom token and a handful of subject photos, applying reconstruction and class-preserving priors during finetuning. While OFT builds upon the DreamBooth paradigm, it diverges by utilizing orthogonal transformations during finetuning, instead of standard parameter updating with reduced learning rates.

\textbf{Controllable generation}. Image-to-image translation represents a subset of controllable generation, where the output is modulated according to specific input conditions. Historically, the predominant approach has involved conditional generative adversarial networks~\cite{isola2017image,zhu2017toward,wang2018high,park2019semantic,choi2018stargan}. The application of diffusion models to this task has also gained traction~\cite{saharia2022palette,voynov2022sketch,wang2022pretraining}. Recently, ControlNet~\cite{zhang2023adding} introduced a mechanism for incorporating extra control signals by finetuning a pretrained diffusion model, thereby achieving notable advances in controllable generation. A related method, T2I-Adapter~\cite{mou2023t2i}, likewise improves the controllability of image generation through targeted finetuning of a diffusion model. Building upon the experimental setup used in \cite{zhang2023adding,mou2023t2i}, we employ OFT to adapt pretrained diffusion models. Our method consistently delivers enhanced controllability while utilizing less training data and fewer parameters during finetuning. Crucially, OFT does not incur any increase in computational cost at inference.

\textbf{Model finetuning}. The paradigm of adapting large pretrained models to specific downstream applications has gained considerable momentum~\cite{devlin2018bert,brown2020language,he2022masked}. Within this broader trend, various adaptation strategies (\emph{e.g.}, \cite{hulora2022,houlsby2019parameter,pfeiffer2020adapterfusion}) have been extensively explored in natural language processing. OFT’s closest analog is LoRA~\cite{hulora2022}, which approaches finetuning as a low-rank, additive modification to the model weights. Distinct from this, OFT implements a layer-wise, shared orthogonal transformation that adjusts neuron weights multiplicatively. This technique theoretically conserves the pairwise angles between neurons in each layer, enhancing model stability during finetuning.

\section{Orthogonal Finetuning}

\subsection{Why Does Orthogonal Transformation Make Sense?}

We begin our exploration by articulating the rationale for employing orthogonal transformations in finetuning text-to-image diffusion models. Our analysis is structured around two fundamental questions: (1) What motivates the finetuning of neuron angles (i.e., their directional component)? (2) What advantages does orthogonal transformation offer for altering these angles during finetuning?

To address the first question, we draw upon empirical findings from \cite{liu2018decoupled,chen2020angular}, which highlight that differences in angular features effectively describe the semantic disparity. SphereNet~\cite{liu2017deep} demonstrates that enforcing normalization of all neuron activations onto a unit hypersphere during network training preserves model capacity and even enhances generalization, indicating that neuronal directions encapsulate the essential data characteristics. To illustrate the significance of neuron angularity, we design a simple experiment depicted in Figure~\ref{fig:toy_ae}. Here, a basic convolutional autoencoder is trained using images of flowers. During training, we rely on the conventional inner product for feature map computation, where $z$ represents the convolutional kernel $\bm{w}$ applied to the input $\bm{x}$ within a sliding window. At test time, we consider three approaches for producing the feature map: (a) the original inner product from training, (b) isolating only the magnitude component, and (c) extracting only the angular component. Results, as shown in Figure~\ref{fig:toy_ae}, reveal that angular information alone enables almost perfect reconstruction of the original images, while magnitude lacks meaningful information. It should be noted that the cosine activation in (c) is not used during training; the training process depends entirely on the inner product. These findings suggest that neuron directions encode the primary semantic information from the images. Thus, adjustments to neuron direction are likely to be most effective for altering image semantics.

With respect to the second question, the most direct method for tuning neuron direction is by collectively applying rotation or reflection to all neurons in a layer, inherently employing orthogonal transformations. Although alternative angular modifications, such as varying rotations for different neurons, could offer greater flexibility, we observe that orthogonal transformations offer an optimal balance between adaptability and structural discipline. Furthermore, it has been established in \cite{liu2021orthogonal} that orthogonal transformations have ample representational power for neural network learning. To reinforce this point, we conduct an experiment to showcase the regularizing benefits provided by enforcing orthogonality. Leveraging the experimental protocol of ControlNet~\cite{zhang2023adding}, we compare outcomes depicted in Figure~\ref{fig:ortho}. The results indicate that our standard OFT maintains consistent performance and achieves precise control post-training (epoch 20), whereas the version of OFT lacking the orthogonality constraint is incapable of generating plausible images or achieving effective control. This experiment underlines the criticality of the orthogonality constraint for OFT.

\subsection{General Framework}

Traditional finetuning generally relies on gradient descent with a reduced learning rate, updating either the entire model or select layers. Employing a small learning rate implicitly restricts changes from the original pretrained weights, and as a result, conventional finetuning effectively aims to minimize the difference $\thickmuskip=2mu \medmuskip=2mu \| \bm{M} - \bm{M}^0\|$ where $\bm{M}$ and $\bm{M}^0$ represent the weights after finetuning and before, respectively. Although this implicit penalty provides an intuitive safeguard against diverging too much from the pretrained parameters, it can permit excessive flexibility—especially problematic when finetuning large-scale models. To counteract this, LoRA incorporates an explicit low-rank restriction on the update, namely $\thickmuskip=2mu \medmuskip=2mu \text{rank}(\bm{M} - \bm{M}^0)=r'$, with $r'$ designated as a modest value. Unlike LoRA, OFT focuses on controlling pairwise neuron similarity, introducing a constraint of the form $\thickmuskip=2mu \medmuskip=2mu \| \text{HE}(\bm{M})-\text{HE}(\bm{M}^0) \|=0$, with $\text{HE}(\cdot)$ denoting the hyperspherical energy associated with a weight matrix. 

To illustrate, consider a fully connected layer expressed as $\thickmuskip=2mu \medmuskip=2mu \bm{W}=\{\bm{w}_1,\cdots,\bm{w}_n\}\in\mathbb{R}^{d\times n}$, where each column $\bm{w}_i\in\mathbb{R}^{d}$ corresponds to the $i$-th neuron, and $\bm{W}^0$ signifies the pretrained parameters. The resulting output $\thickmuskip=2mu \medmuskip=2mu \bm{z}\in\mathbb{R}^n$ for this layer is given by $\thickmuskip=2mu \medmuskip=2mu \bm{z}=\bm{W}^\top\bm{x}$ for input $\thickmuskip=2mu \medmuskip=2mu \bm{x}\in\mathbb{R}^d$. The core motivation of OFT is to minimize alterations in hyperspherical energy between finetuned and original models, formalized as

\begin{equation}\label{eq:oft_principle}
\footnotesize
\min_{\bm{W}} \left\| \text{HE}(\bm{W})-\text{HE}(\bm{W}^0)\right\|~~\Leftrightarrow~~\min_{\bm{W}} \bigg{\|} \sum_{i\neq j} \|\hat{\bm{w}}_i-\hat{\bm{w}}_j\|^{-1}-\sum_{i\neq j} \|\hat{\bm{w}}^0_i-\hat{\bm{w}}^0_j\|^{-1}\bigg{\|}
\end{equation}

where $\thickmuskip=2mu \medmuskip=2mu \hat{\bm{w}}_i:=\bm{w}_i/\|\bm{w}_i\|$ denotes the normalized $i$-th neuron, and $\text{HE}(\bm{W}):= \sum_{i\neq j} \|\hat{\bm{w}}_i-\hat{\bm{w}}_j\|^{-1}$ characterizes the hyperspherical energy for $\bm{W}$. It is evident that Eq.~\eqref{eq:oft_principle} achieves its minimum—zero—when $\bm{W}$ and $\bm{W}^0$ are related via orthogonal transformations; in other words, when $\bm{W}=\bm{R}\bm{W}^0$ for some orthogonal matrix $\thickmuskip=2mu \medmuskip=2mu \bm{R}\in\mathbb{R}^{d\times d}$ (where determinant $1$ indicates rotation and $-1$ reflection). OFT leverages this property by proposing that only orthogonal transformations—using layer-wise shared orthogonal matrices—are needed for effective adaptation. Thus, given a pretrained weight matrix $\thickmuskip=2mu \medmuskip=2mu\bm{W}^0\in\mathbb{

\begin{equation}\label{eq:oft_general}
    \footnotesize
    \bm{z}=\bm{W}^\top\bm{x}=(\bm{R}\cdot\bm{W}^0)^\top\bm{x},~~~\text{s.t.}~\bm{R}^\top\bm{R}=\bm{R}\bm{R}^\top=\bm{I}
\end{equation}

In this formulation, $\bm{W}$ refers to the weight matrix that has been adapted via OFT finetuning, while $\bm{I}$ stands for the identity matrix. An overview of OFT can be seen in Figure~\ref{fig:block}. Analogous to the zero initialization technique applied in LoRA, OFT requires the finetuning procedure to commence precisely from the pretrained model parameters $\bm{W}^0$. To ensure this starting point, we set the initial orthogonal matrix $\bm{R}$ to be the identity, thereby aligning the finetuned parameters initially with the pretrained values. Maintaining the orthogonality of $\bm{R}$ is essential, which can be accomplished using differentiable orthogonalization approaches described in \cite{liu2021orthogonal,lezcano2019cheap}. We elaborate below on efficient techniques to enforce orthogonality in practice.

\subsection{Efficient Orthogonal Parameterization}\label{sect:eff_ortho}

Conventional orthogonalization techniques, such as the Gram-Schmidt process, are theoretically differentiable but are prohibitively expensive to compute for large-scale matrices~\cite{liu2021orthogonal}. To improve efficiency, we employ the Cayley parameterization to construct the orthogonal matrix. More concretely, the orthogonal matrix is represented as $\thickmuskip=2mu \medmuskip=2mu \bm{R}=(\bm{I}+\bm{Q})(I-\bm{Q})^{-1}$, where $\bm{Q}$ is a skew-symmetric matrix that fulfills $\thickmuskip=2mu \medmuskip=2mu \bm{Q}=-\bm{Q}^\top$. This parameterization greatly enhances computational efficiency, although it imposes a minor constraint—namely, Cayley parameterization only generates orthogonal matrices with determinant $1$, restricting them to the special orthogonal group. Our empirical results indicate that this restriction does not negatively impact performance. Despite adopting the Cayley transform, when the dimensionality $d$ is high, $\bm{R}$ can still involve a substantial number of parameters, compromising efficiency. To overcome this, we introduce a block-diagonal structure for $\bm{R}$, partitioning it into $r$ blocks. The resultant structure is given by:

\begin{equation}
    \footnotesize
    \bm{R}=\text{diag}(\bm{R}_1,\bm{R}_2,\cdots,\bm{R}_r)=\begin{bmatrix}
    \bm{R}_1\in \text{O}(\frac{d}{r}) &  &\\
    &\ddots & \\
    & & \bm{R}_r\in \text{O}(\frac{d}{r})
    \end{bmatrix}\in \text{O}(d)
\end{equation}

Here, $\text{O}(d)$ signifies the orthogonal group in $d$ dimensions, where both $\thickmuskip=2mu \medmuskip=2mu \bm{R}\in\mathbb{R}^{d\times d}$ and $\thickmuskip=2mu \medmuskip=2mu \bm{R}_i\in\mathbb{R}^{d/r\times d/r},\forall i$ represent orthogonal matrices. Setting $\thickmuskip=2mu \medmuskip=2mu r=1$ results in a conventional unconstrained orthogonal matrix, as the block-diagonal structure reduces to a standard form. For an orthogonal matrix of shape $\thickmuskip=2mu \medmuskip=2mu d\times d$, there are $\thickmuskip=2mu \medmuskip=2mu d(d-1)/2$ free parameters, corresponding to $\mathcal{O}(d^2)$ computational complexity. In contrast, an $r$-block diagonal orthogonal matrix requires only $\thickmuskip=2mu \medmuskip=2mu d(d/r-1)/2$ parameters, leading to a reduced complexity of $\thickmuskip=2mu \medmuskip=2mu \mathcal{O}(d^2/r)$. By choosing to reuse (or share) block matrices, that is, setting $\thickmuskip=2mu \medmuskip=2mu\bm{R}_i=\bm{R}_j,\forall i\neq j$, one can decrease the parameter count even further, down to $\mathcal{O}(d^2/r^2)$. Importantly, all these efficiency-improving approaches still preserve the orthogonality of $\bm{R}$, ensuring that the hyperspherical energy remains intact.

A comparison of OFT and LoRA in terms of parameter count shows significant differences in efficiency. For LoRA with low-rank parameter $r'$, the count of additional trainable parameters is $\thickmuskip=2mu \medmuskip=2mu r'(d+n)$. If both $r$ and $r'$ are tied to the feature dimension $d$ (\emph{e.g.}, $\thickmuskip=2mu \medmuskip=2mu r=r'=\alpha d$ with $\thickmuskip=2mu \medmuskip=2mu 0<\alpha\leq 1$ as some scalar), then the parameter requirement for LoRA is $\mathcal{O}(d^2+dn)$, whereas for OFT it becomes $\mathcal{O}(d)$. To concretize this difference, consider a $128\times128$ weight matrix. For $\thickmuskip=2mu \medmuskip=2mu r'=8$, LoRA introduces $2,048$ parameters, whereas OFT with $\thickmuskip=2mu \medmuskip=2mu r=8$ (without any block sharing) utilizes $960$ trainable parameters.

\subsection{Constrained Orthogonal Finetuning}

The flexibility of the OFT method can be further reduced by ensuring that the finetuned model stays within a limited region surrounding the initial pretrained weights. The COFT approach accomplishes this by employing the following computation for each forward pass:

\begin{equation}\label{eq:coft_general}
    \footnotesize
    \bm{z}=\bm{W}^\top\bm{x}=(\bm{R}\cdot\bm{W}^0)^\top\bm{x},~~~\text{s.t.}~\bm{R}^\top\bm{R}=\bm{R}\bm{R}^\top=\bm{I},~\norm{\bm{R}-\bm{I}}\leq \epsilon
\end{equation}

This formulation applies both an orthogonality requirement and an $\epsilon$-ball constraint around the identity. Section~\ref{sect:eff_ortho} presents the Cayley parameterization, which enforces orthogonality. Nevertheless, directly introducing the $\epsilon$-deviation condition to a Cayley-parameterized matrix is challenging. To better understand the Cayley transform, the Neumann series expansion is invoked to approximate $\thickmuskip=2mu \medmuskip=2mu \bm{R}=(\bm{I}+\bm{Q})(I-\bm{Q})^{-1}$ as $\thickmuskip=2mu \medmuskip=2mu

The convergence of the series in the operator norm allows us to internalize the constraint $\thickmuskip=2mu \medmuskip=2mu\|\bm{R}-\bm{I}\|\leq\epsilon$ within the Cayley transform. As a result, an equivalent condition emerges: $\thickmuskip=2mu \medmuskip=2mu \|\bm{Q}-\bm{0}\|\leq\epsilon'$, where $\bm{0}$ represents the zero matrix and $\epsilon'$ denotes a distinct error tolerance parameter from $\epsilon$. This revised condition on $\bm{Q}$ lends itself well to straightforward enforcement through projected gradient descent. To ensure identity initialization for the orthogonal transformation $\bm{R}$, we set the initial value of $\bm{Q}$ to the zero matrix. In essence, COFT incorporates two explicit forms of regularization: limiting the hyperspherical energy discrepancy, and explicitly constraining deviations from the pretrained weights. While traditional finetuning techniques often adopt the second constraint implicitly, COFT articulates it as a formal restriction. In practice, although OFT performs robustly, COFT demonstrates improved finetuning stability, likely attributable to its direct constraint on deviation. Figure~\ref{fig:eps_coft} showcases how varying $\epsilon$ influences COFT's effectiveness: as $\epsilon$ grows, the output of COFT aligns more closely with that of OFT, while smaller values of $\epsilon$ lead COFT to preserve greater similarity to the original pretrained text-to-image diffusion network.

\subsection{Re-scaled Orthogonal Finetuning}

We introduce an augmentation to the standard OFT method by learning an additional scalar scaling parameter for each neuron. The rationale for this design is that rescaling neurons does not impact the hyperspherical energy, given that normalization nullifies the magnitude. The forward operation becomes: $\thickmuskip=2mu \medmuskip=2mu\bm{z}=(\bm{R}\bm{W}^0\bm{D})^\top\bm{x}$\footnote{\textbf{Errata}: The NeurIPS camera-ready paper incorrectly presents the re-scaled OFT forward pass as $\thickmuskip=2mu \medmuskip=2mu\bm{z}=(\bm{D}\bm{R}\bm{W}^0)^\top\bm{x}$. However, our original code is accurate, and thus the results in Appendix~\ref{app:rescaled} remain valid.}, where $\thickmuskip=2mu \medmuskip=2mu \bm{D}=\text{diag}(s_1,\cdots,s_n)\in\mathbb{R}^{n\times n}$ stands for a diagonal matrix with positive learnable entries $s_1,\cdots,s_n$. This approach contrasts with OFT's conventional forward computation in Eq.~\eqref{eq:oft_general}, which restricts learnable parameters to $\bm{R}$. Here, both $\bm{D}$ and $\bm{R}$ are trainable. The enhanced re-scaled OFT increases expressivity at the cost of only a minimal number of extra parameters. In our main experiments, we retain the original OFT to highlight the capability of orthogonal transformations, but we observe that re-scaled OFT generally exhibits superior results (refer to Appendix~\ref{app:rescaled}).

\section{Intriguing Insights and Discussions}

\textbf{OFT's architectural versatility}. Fundamentally, OFT can be integrated into any neural network architecture. For instance, in the context of Transformers, LoRA typically modifies the attention weight matrices~\cite{hulora2022}. To ensure comparability with LoRA, our experiments restrict OFT to the finetuning of attention parameters. Beyond standard fully connected layers, OFT is inherently compatible with convolutional layers, as its block-diagonal $\bm{R}$ structure admits useful interpretations for convolutions (a flexibility not afforded by LoRA). By selecting the number of blocks in $\bm{R}$ to match the quantity of input channels, we attain block-wise transformations where each block operates exclusively on a particular neuron channel, akin to adapting depth-wise convolution filters~\cite{chollet2017xception}. If, instead, all blocks of $\bm{R}$ are tied, then OFT applies a uniform orthogonal transformation across all channels. Our preliminary experiments on employing OFT for convolutional layer finetuning are detailed in Appendix~\ref{app:convolution}.

\textbf{Connection to LoRA}. LoRA achieves parameter efficiency by introducing a low-rank matrix, thereby minimizing large alterations in the original pretrained weight matrix. By contrast, OFT employs an orthogonal (thus full-rank) transformation on the pretrained weights, constraining the transformation such that it does not compromise the knowledge already encoded. The forward computation in OFT can be formulated as $\thickmuskip=2mu \medmuskip=2mu \bm{z}=(\bm{R}\bm{W}^0)^\top\bm{x}=(\bm{W}^0+(\bm{R}-\bm{I})\bm{W}^0)^\top\bm{x}$, where the term $\thickmuskip=2mu \medmuskip=2mu (\bm{R}-\bm{I})\bm{W}^0$ serves a functionally similar role to LoRA's low-rank adaptation. If $\bm{W}^0$ has full rank (as is commonly the case), then whenever $\thickmuskip=2mu \medmuskip=2mu \bm{R}-\bm{I}$ is of low rank, OFT can likewise be seen as carrying out a low-rank modification of the weights. Both methods utilize hyperparameters—rank $r'$ in LoRA and diagonal block size $r$ in OFT—to limit the number of trainable parameters. Fundamentally, LoRA and OFT exemplify distinct strategies for parameter efficiency: LoRA leverages low-rank structures, while OFT capitalizes on sparsity induced by block-diagonal orthogonality constraints.

\textbf{Why OFT converges faster}. As indicated in Figure~\ref{fig:toy_ae}, making adjustments to neuron orientations is highly effective for altering representations—precisely the mechanism OFT is engineered to implement. Additionally, OFT’s update process can be interpreted as optimization on a smooth hypersphere manifold, which improves the optimization landscape and accelerates convergence, an advantage also reported in \cite{liu2021orthogonal}.

\textbf{Why not minimize hyperspherical energy}. In contrast to \cite{liu2021orthogonal}, our approach does not pursue minimization of hyperspherical energy. For classification settings, it is preferable for neurons to be non-redundant; minimizing hyperspherical energy results in an even distribution of neurons on the sphere, which is not necessarily advantageous for finetuning and can inadvertently degrade valuable information learned during pretraining.

\textbf{Balancing flexibility and regularization in finetuning}. Our analysis reveals an inherent tension between model adaptability and imposed constraints. While traditional finetuning achieves maximal adaptability, it is prone to instability and often leads to model degeneration. The OFT method, despite its simplicity, adeptly negotiates this balance by maintaining invariant pairwise angles between neurons. Additionally, employing block-diagonal parameterization enforces a stricter form of regularization on the orthogonal transformation.

\textbf{Inference remains efficient}. In contrast to ControlNet, the OFT framework does not increase inference time for the adapted model. At inference, the orthogonal matrix $\bm{R}$ can be directly merged into the pretrained weights $\bm{W}^0$, yielding an equivalent parameter matrix $\thickmuskip=2mu \medmuskip=2mu \bm{W}=\bm{R}\bm{W}^0$. Consequently, the inference performance matches that of the original pretrained system.

\section{Experiments and Results}

\textbf{Experimental configuration}. The experimental evaluation is conducted using Stable Diffusion v1.5~\cite{rombach2022high} as the pretrained text-to-image backbone. Images for each method are selected at random for objective comparison. Subject-driven generation experiments largely adhere to the DreamBooth~\cite{ruiz2023dreambooth} protocol, while experiments for controllable generation are consistent with the methodologies described in ControlNet~\cite{zhang2023adding} and T2I-Adapter~\cite{mou2023t2i}. In comparison to LoRA, OFT and COFT are applied strictly to the same network layers as LoRA. Further experimental details and additional results are presented in Appendix~\ref{app:settings}.

\subsection{Subject-driven Generation}

\textbf{Experimental details}. DreamBooth~\cite{ruiz2023dreambooth} and LoRA~\cite{hulora2022} serve as comparative baselines in our experiments. All approaches optimize the loss function prescribed by DreamBooth. For both DreamBooth and LoRA, procedures and hyperparameters generally adhere to those recommended in their respective papers. Supplementary results can be found in Appendix~\ref{app:settings}, \ref{app:coft_oft}, \ref{app:more_qualitative}, \ref{app:failure}.

\textbf{Finetuning stability and convergence}. We begin by assessing the stability during finetuning and the convergence rates for DreamBooth, LoRA, OFT, and COFT. The outcomes are illustrated in Figure~\ref{fig:teaser} and Figure~\ref{fig:db_convergence}. Notably, COFT and OFT exhibit stable finetuning behavior for the diffusion model. By the 400th iteration, both DreamBooth and the OFT-based variants provide effective subject control, whereas LoRA struggles to maintain subject identity. When training reaches 2000 iterations, DreamBooth demonstrates mode collapse in its outputs, and LoRA is unable to correctly depict a yellow shirt—often producing yellow fur instead. In sharp contrast, OFT and COFT continue to provide stable and reliable control over image generation throughout. These findings demonstrate that the OFT framework converges both quickly and robustly for subject-driven tasks. Additionally, the reduced sensitivity to the number of finetuning steps is highly advantageous, as it minimizes the need to manually adjust iteration counts for varying subjects. With both OFT and COFT, it is feasible to predefine a sufficiently large number of iterations without intricate hyperparameter tuning. For COFT, an appropriate choice of $\epsilon$ makes both the learning rate and iteration count essentially trivial to configure.

\textbf{Quantitative comparison}. Adopting the evaluation protocols from \cite{ruiz2023dreambooth}, we perform a quantitative analysis for subject fidelity (DINO~\cite{caron2021emerging}, CLIP-I~\cite{radford2021learning}), prompt alignment (CLIP-T~\cite{radford2021learning}), and sample diversity (LPIPS~\cite{zhang2018unreasonable}). Here, CLIP-I reflects the mean pairwise cosine similarity between the CLIP embeddings of generated and corresponding real images. The DINO score is computed in an analogous manner using ViT S/16 DINO features. For CLIP-T, the metric captures the mean cosine similarity between CLIP embeddings of the generated images and the conditioning text. Furthermore, we utilize LPIPS to assess the mean cosine similarity across samples depicting the same subject conditioned on identical prompts. According to Table~\ref{table:dreambooth_final}, both COFT and OFT deliver substantial improvements over DreamBooth and LoRA in DINO and CLIP-I subject fidelity measures, while maintaining slightly better or equivalent results on text prompt and diversity metrics. For all approaches, the experiments are repeated 30 times with distinct random seeds to reduce variance. Altogether, these results demonstrate that the OFT framework excels not only in achieving superior stability and convergence but also in consistently delivering enhanced final performance.

\textbf{Qualitative comparison}. To provide clearer insights into the advantages of OFT, we present a selection of randomly chosen instances demonstrating subject-driven generation in Figure~\ref{fig:dreambooth_final}. For consistency and fairness, each example is produced using the same model finetuned for each respective approach, ensuring that no individual text prompt is optimized for any specific output. For every compared method, we select the model exhibiting the highest validation CLIP scores. Analyzing the visual outcomes in Figure~\ref{fig:dreambooth_final}, we note that both OFT and COFT maintain robust semantic subject fidelity, in contrast to LoRA, which frequently struggles to retain subject characteristics (for instance, in the bowl example, LoRA entirely loses the subject identity). Furthermore, OFT and COFT provide superior precision in controlling the output via text prompts, whereas DreamBooth, despite excelling at subject identity preservation, often fails to align generated content with textual input (such as in the bowl example’s first row). These qualitative results substantiate that the OFT framework achieves a stronger balance of controllability and subject preservation. Additionally, the number of iterations in OFT does not have a significant impact on its performance, enabling OFT to maintain high quality even with a substantial number of iterations—a capability not observed in either DreamBooth or LoRA. Additional qualitative results can be found in Appendix~\ref{app:more_qualitative}. Human assessment experiments reported in Appendix~\ref{app:human} further corroborate the superior performance of OFT.

\subsection{Controllable Generation}

\textbf{Settings}. We evaluate OFT against several established baselines: ControlNet~\cite{zhang2023adding}, T2I-Adapter~\cite{mou2023t2i}, and LoRA~\cite{hulora2022}. Our experiments focus on three challenging controllable generation tasks highlighted in the main text: transforming Canny edges to images~(C2I) using the COCO dataset~\cite{lin2014microsoft}, converting segmentation maps to images~(S2I) on the ADE20K dataset~\cite{zhou2017scene}, and generating faces from landmarks~(L2F) on the CelebA-HQ dataset~\cite{karras2018progressive,xia2021tedigan}. Every approach involves finetuning the Stable Diffusion~(SD) v1.5 model across all three datasets for 20 training epochs. Expanded results and further examples can be found in Appendix~\ref{app:more_qualitative},~\ref{app:more_control_task},~\ref{app:failure}.

\textbf{Convergence}. We assess how rapidly ControlNet, T2I-Adapter, LoRA, and COFT converge when applied to the L2F task, utilizing both numerical metrics and visual inspection. For our quantitative analysis, we calculate the mean $\ell_2$ distance between the predicted face landmarks and the control landmarks. Figure~\ref{fig:face_convergence} presents the landmark error trajectories as models are finetuned across varying numbers of epochs. Notably, COFT and OFT display a much quicker convergence rate compared to baselines; for instance, LoRA requires 20 epochs to reach the accuracy that OFT obtains after only 8 epochs. Importantly, both OFT and COFT have a similar count of trainable parameters to LoRA—substantially less than ControlNet—yet they are considerably more efficient in training convergence. The efficiency of OFT’s convergence is further corroborated by the evidence shown in Figure~\ref{fig:teaser}, where OFT outperforms both ControlNet and LoRA in terms of data efficiency, achieving satisfactory convergence with only 5\% of the ADE20K dataset. In qualitative comparisons, our analysis focuses on OFT, COFT, and ControlNet, as ControlNet yields landmark errors most similar to ours. Visualizations in Figure~\ref{fig:face_convergence_qual} illustrate that OFT and COFT consistently exhibit stable convergence, with the synthesized face poses progressively conforming to the control landmarks. Conversely, ControlNet demonstrates an abrupt onset of controllability after the 8th epoch, mirroring the sudden convergence phenomenon reported by \cite{zhang2023adding}. This enhanced stability in OFT’s convergence behavior contributes to its practical utility: the predictable and gradual learning dynamics facilitate the search for optimal hyperparameters, making OFT straightforward to deploy in real-world scenarios.

\textbf{Quantitative comparison}. To assess the effectiveness of controllable generation, we propose a control consistency metric. This metric involves deriving the control signal from the synthesized image and measuring its correspondence to the initial control signal provided as input. Specifically, for the C2I task, we utilize IoU and F1 score, while for the S2I task, we report mean IoU along with mean and overall accuracy. The L2F task is evaluated based on the average $\ell_2$ distance between ground truth control landmarks and the predicted counterparts. Additional information about these consistency measures can be found in Appendix~\ref{app:settings}. Each competing method is evaluated using optimal hyperparameter choices. Table~\ref{table:control_gen} presents the results, demonstrating that both OFT and COFT achieve significantly higher control fidelity and precision relative to alternative techniques. It is also evident that methods relying on adapters (such as T2I-Adapter and ControlNet) exhibit slower convergence and attain subpar end performance. While LoRA surpasses ControlNet in S2I tasks, it underperforms on C2I and L2F tasks. Overall, our observations indicate that the achievable upper bound for LoRA is somewhat limited, despite extensive hyperparameter optimization. In contrast, our empirical studies suggest that the performance of the OFT framework continues to improve as the number of trainable parameters is increased, indicating no evident saturation point. We would like to highlight that our approach to quantitative evaluation for controllable generation represents a novel contribution, providing reliable assessment of model controllability that aligns closely with the capabilities of existing segmentation or detection systems.

\textbf{Qualitative comparison}. We further conduct a qualitative analysis of OFT and COFT in comparison to contemporary leading methods such as ControlNet, T2I-Adapter, and LoRA. As illustrated by the randomly sampled outputs in Figure~\ref{fig:control_final}, OFT and COFT consistently deliver both high-fidelity, realistic imagery and precise controllability. In the S2I scenario, LoRA fails to produce outputs that correspond to the input segmentation maps, whereas ControlNet, OFT, and COFT maintain strong control over generated content. Distinct from ControlNet, OFT and COFT achieve images of superior fidelity, enriched detail, and more plausible geometrical structure, all while requiring significantly fewer parameters. Regarding the C2I task, OFT and COFT demonstrate the ability to imaginatively synthesize semantically coherent images from basic Canny edge inputs, a capacity where T2I-Adapter and LoRA underperform. In the L2F setting, our approach yields the most precise facial pose alignment, even with challenging orientations. Across the three evaluated control tasks, the results indicate that OFT and COFT surpass the baseline state-of-the-art models in image quality, underscoring the efficacy of the OFT framework for controlled image generation. To ensure a thorough qualitative assessment, additional examples for each task can be found in Appendix~\ref{app:control_exp}, while Appendix~\ref{app:more_control_task} extends this evaluation by showcasing OFT’s strong performance across diverse control tasks such as dense pose to human body, sketch-to-image, and depth-to-image translation.

\section{Concluding Remarks and Open Problems}

Inspired by the insight that the angular relationships among neurons play a pivotal role in defining visual semantics, this work introduces a straightforward yet robust finetuning approach—orthogonal finetuning—for the manipulation of text-to-image diffusion models. Our method focuses on two main applications: subject-driven generation and controllable generation. Relative to prior approaches, OFT achieves enhanced control and stable finetuning while utilizing a smaller number of parameters. Crucially, OFT adds no extra computational burden during inference, making it suitable for efficient deployment.

The proposed OFT method also raises several promising avenues for future research. Firstly, enforcing orthogonality in OFT relies on Cayley parametrization, which requires a matrix inversion and thereby imposes certain scalability constraints. Although our adoption of block diagonal parametrization partly addresses this, finding differentiable and more efficient methods for computing the matrix inverse remains an open challenge. Secondly, OFT offers distinctive opportunities for compositionality: orthogonal matrices obtained from separate OFT finetuning tasks can be multiplied to yield another orthogonal matrix. Investigating whether such a composition retains the task-specific knowledge across multiple downstream applications is an intriguing topic for future exploration. Lastly, the parameter efficiency of OFT largely hinges on the use of a block diagonal structure, which can introduce additional biases and restrict flexibility. Identifying strategies to further enhance parameter efficiency while minimizing bias continues to be an essential open problem.

\end{document}